
\vspace{-0.8cm}

This chapter presents the system analysis and design of the proposed \textbf{Automated Car Parking System}. It covers the main requirements, system architecture, embedded system diagram, hardware setup, circuit connection, robotic arm gate structure, and design rationale.

\section{Requirement Analysis}

Requirement analysis defines what the system should perform and how it should behave. The requirements were identified by observing common parking problems, reviewing existing smart parking systems, and discussing the expected system behavior among group members. The main stakeholders are drivers, parking administrators, and system maintainers.

The system must detect vehicles, verify authorized users, operate the robotic arm gate, update parking slot information, display real-time status, and send parking records to an online server or database. The requirements are divided into functional and non-functional requirements.

\subsection{Functional Requirements}

\begin{table}[H]
\centering
\small
\caption{Functional Requirements}
\label{tab:functional_requirements}
\begin{tabular}{p{2.2cm} p{11cm}}
\toprule
\textbf{ID} & \textbf{Requirement Description} \\
\midrule
FR1 & Detect vehicles at the entrance using an IR sensor. \\
FR2 & Detect vehicles at the exit using an IR sensor. \\
FR3 & Verify authorized users using RFID cards. \\
FR4 & Open and close the two-link robotic arm gate using a servo motor. \\
FR5 & Update available parking slot count after valid entry and exit. \\
FR6 & Display current slot availability on a 16x2 I2C LCD. \\
FR7 & Send entry and exit records to an online server using Wi-Fi. \\
FR8 & Prevent vehicle entry when no parking slot is available. \\
\bottomrule
\end{tabular}
\normalsize
\end{table}

\subsection{Non-Functional Requirements}

\begin{table}[H]
\centering
\small
\caption{Non-Functional Requirements}
\label{tab:nonfunctional_requirements}
\begin{tabular}{p{2.2cm} p{11cm}}
\toprule
\textbf{ID} & \textbf{Requirement Description} \\
\midrule
NFR1 & The system should respond quickly after vehicle detection. \\
NFR2 & The robotic arm gate should move smoothly and safely. \\
NFR3 & The system should be simple and user-friendly. \\
NFR4 & The design should be low-cost and suitable for prototype implementation. \\
NFR5 & The slot count should remain accurate during entry and exit. \\
NFR6 & Online records should be transmitted reliably when Wi-Fi is available. \\
\bottomrule
\end{tabular}
\normalsize
\end{table}

\section{System Overview / Architecture}

The Automated Car Parking System is an embedded and IoT-based system. It uses IR sensors for vehicle detection, RFID for authentication, Arduino UNO for processing, a servo motor-controlled two-link robotic arm for gate movement, LCD for local display, and ESP-01S Wi-Fi module for online data transmission.

During entry, the system detects a vehicle, checks slot availability, verifies RFID, opens the robotic arm gate, updates the slot count, and sends the entry record online. During exit, the system detects the vehicle, opens the gate, updates the slot count, and sends the exit record.

\subsection{High-Level Architecture}

Fig.~\ref{fig:system_workflow} shows the high-level workflow of the system. The Arduino UNO acts as the central controller and communicates with the sensors, RFID module, LCD, servo motor, and Wi-Fi module.

\begin{figure}[H]
    \centering
    \includegraphics[width=0.78\textwidth]{system_workflow.png}
    \caption{High-Level Workflow of Automated Car Parking System}
    \label{fig:system_workflow}
\end{figure}

\subsection{Embedded System Diagram}

The embedded system diagram in Fig.~\ref{fig:embedded_system_diagram} shows the interaction among the hardware and software modules. The Arduino UNO receives input from the entry IR sensor, exit IR sensor, and RFID reader. It then controls the servo motor-based two-link robotic arm gate, updates the LCD, and communicates with the ESP-01S Wi-Fi module.

The ESP-01S sends parking event data to the middleware server, which stores the information in the online database. The stored data can later be monitored through an online interface or mobile application.

\begin{figure}[H]
    \centering
    \includegraphics[width=0.82\textwidth]{embedded_car_parking_system_diagram.png}
    \caption{Embedded System Diagram of Automated Car Parking System}
    \label{fig:embedded_system_diagram}
\end{figure}

\section{Hardware Setup}

The hardware setup includes the main electronic and mechanical components required for vehicle detection, authentication, gate control, display, and communication.

\begin{table}[H]
\centering
\small
\caption{Main Hardware Components}
\label{tab:hardware_components}
\begin{tabular}{p{4cm} p{9cm}}
\toprule
\textbf{Component} & \textbf{Function} \\
\midrule
Arduino UNO & Main controller of the system. \\
IR Sensors & Detect vehicle presence at entry and exit points. \\
RFID Module & Reads RFID cards for user authentication. \\
Servo Motor & Controls the two-link robotic arm gate. \\
Two-Link Robotic Arm & Works as the automatic gate mechanism. \\
16x2 I2C LCD & Displays slot availability and system messages. \\
ESP-01S Wi-Fi Module & Sends parking data to the online server or database. \\
Power Supply & Provides required power to the circuit and components. \\
\bottomrule
\end{tabular}
\normalsize
\end{table}

\subsection{Circuit Connection}

The circuit connects all modules with the Arduino UNO. The IR sensors are connected to digital pins for vehicle detection. The RFID module uses SPI communication. The LCD uses I2C communication, which reduces wiring complexity. The servo motor receives control signals from the Arduino, and the ESP-01S communicates through serial AT commands for internet connectivity.

\begin{figure}[H]
    \centering
    \includegraphics[width=0.78\textwidth]{circuit_connection.png}
    \caption{Circuit Connection of the Proposed System}
    \label{fig:circuit_connection}
\end{figure}

\subsection{Robot Structure Details}

The gate mechanism is designed using a servo motor-controlled two-link robotic arm. The two-link structure provides controlled robotic movement instead of a simple static barrier. When the system grants permission, the servo motor moves the arm to the open position. After the vehicle passes, the arm returns to the closed position.

This structure makes the system suitable for a robotics lab project because it demonstrates actuator control, mechanical movement, and automated response based on sensor input.

\section{Design Alternatives and Rationale}

Several alternatives were considered during design. A manual gate was rejected because it does not provide automation. A simple servo barrier was easier to build but did not demonstrate enough robotic motion. A camera-based system was also considered but would increase cost and complexity.

The final design was selected because it balances cost, functionality, and simplicity. Arduino UNO, IR sensors, RFID, LCD, ESP-01S, and a two-link robotic arm provide a practical solution for a low-cost robotics and IoT-based parking prototype. The design also allows future improvements such as mobile app integration, camera-based vehicle recognition, or payment support.

\section{Summary}

This chapter described the system analysis and design of the Automated Car Parking System. The functional and non-functional requirements, system architecture, embedded system diagram, hardware setup, circuit connection, robotic arm gate mechanism, and design rationale were presented. The selected design is simple, low-cost, and suitable for demonstrating robotics and IoT-based parking automation.

